\scanpage{39}

\uline{L'idée de W.M.}

\ldots\ldots

Donnons quelques détails sur ce que font \uline{réellement} W.M. Faute de pouvoir \unsure{relever} en général un $X_0$ lisse quelconque, on travaill. \inserted{\uline{d'abord}} avec des $X_0$ \uline{affines}. On sait (SGA 1960/61, $\overline{\underline{\mathrm{III}}}$) que un tel $X_0$ se relève de proche en proche en un $\mathfrak X$ \uline{formel} \inserted{\uline{lisse}} sur $\mathfrak S=\operatorname{Spf}W(k)$, \inserted{unique à isom (\uline{non unique}) près,} la première idée serait de prendre $\Omega^\bullet_{\mathfrak X/\mathfrak S}$ défini par passage à la limite à partir des $\Omega^\bullet_{X_n/W_n}$, et de poser
\[
H(X_{\unsure{0|m}})=\mathbb H^*(\mathfrak X,\Omega^\bullet_{\mathfrak X/\mathfrak S})
=H^*(\Gamma(\mathfrak X,\Omega^\bullet_{\mathfrak X/S})\unsure{)}
\]

Mais montrons que cela ne donne pas du tout ce qu'il faut, $\overline{\mathrm m}$ si $X_0=\operatorname{Spec} k[T]$ est la droite affine, donc
\[
\mathfrak X=\operatorname{Spf}
\underbrace{W\{t\}}_{\varprojlim W_n[T]}
\]
\inserted{séries formelles restreintes, i.e. à coeff. $\longrightarrow 0$.}
On voudrait trouver en effet que
\[
\mathbb H^1(\mathfrak X,\Omega^\bullet_{\mathfrak X/S})=0
\]

\margin{38}
